\section{Introduction}
\label{sec:intro}

An agent that trades against rivals faces three problems in sequence. It must judge a deal by what the deal gives the other side as well as by what it gets. It must turn what it acquires into income, which takes cash, a legal opportunity and a correctly specified action. And it must stay solvent, because a debt it cannot pay hands everything it built to someone else. Evaluations of language-model agents usually examine one of these problems at a time, in a bargaining game, a business simulation or a planning task.

Monopoly puts all three on one path. A player completes a color group, often by trade. Houses on the group raise its rent many times over. A rival who lands there and cannot pay goes bankrupt, and in the implementation studied here the creditor inherits the debtor's deeds with their buildings intact. This paper asks where along that path twenty games were decided, and what four LLM systems did and said at each step.

The data are twenty complete four-player games among four configured systems, which we call Sol, Gemini, Opus and Grok after their underlying models. The games contain 6,989 decisions and all end by bankruptcy without a turn limit. For every decision the record holds the legal menu, each response attempt, the action the engine applied and the events that followed. We reconstruct all 526 negotiation threads, all 209 episodes in which a player came to own a complete color group, all 60 eliminations and every transfer of cash and deeds. Three findings organize the paper.

\emph{Bargaining is priced and sensitive to the rival's gain} (Section~\ref{sec:bargaining}). More than half of all offers would complete a color group for someone, and the side that would complete is nearly always the side paying. An offer that completes only its writer's group is rejected outright more than half the time and accepted about one time in fifteen. An offer that completes only the recipient's group is usually answered with a counter. The four systems accept at similar frequencies and decline in different ways, and completing deeds are released mostly by holders who are nearly out of cash.

\emph{A completed group is not yet a developed one} (Section~\ref{sec:development}). The four systems complete similar numbers of groups by trade. Two build on nearly every one at the first opportunity. The other two build on fewer than half, and their groups collect about a fifth as much rent. Most of the unbuilt groups are not decisions to wait. They are purchases that arrived mortgaged or left no cash, engine defaults after invalid build plans, most of them following a refused first attempt, and responses whose text named one action while the tool call named another.

\emph{Bankruptcy moves more property than trade} (Section~\ref{sec:insolvency}). Fifty-six of sixty eliminations could not have been paid from the debtor's own assets even on a generous bound. Estates complete almost twice as many color groups as trades do, and the eventual winner receives the first estate in eighteen games. The systems bargain over this. They sell a deed as insurance against a rival's inheritance, pay above their own stated ceiling to keep a deed out of an estate, and withdraw an offer once its recipient is nearly broke.

Seven incidents with verbatim excerpts accompany the counts. The task permitted public bluffing and deception, so a misleading public statement is bargaining conduct here and is not a breach of instructions.

These findings matter for anyone who expects language-model agents to hold and trade assets. The systems did not close deals readily. They declined or countered 88 percent of offers, often with stated reasons about a rival's position, though frequencies alone cannot show whether they declined the right ones. The value of a negotiated asset was then set by the execution layer as much as by the bargain, and one configuration lost most of that value to an interaction between a provider refusal and an action format. And when default transfers assets to the creditor, a counterparty's solvency becomes something to bargain over.

This is an observational study of four complete configurations, including their prompts, memory, retry policies and provider behavior. Two configurations ran under conditions that change what their records mean, and Section~\ref{sec:setup} states those conditions before any result. The counts are censuses of a finite cohort with strong dependence within games and seed blocks. We do not rank models, estimate the effect of a rival's completion on refusal, or claim that any trade, purchase or bankruptcy was optimal. Section~\ref{sec:discussion} maps the evidence onto eight broader questions about strategic agents and says which ones this cohort cannot answer.

The work builds on economic and bargaining evaluations that record process as well as outcome~\cite{backlund2025vendingbench,li2026arena,wu2026agentmonopoly,hu2026finboardbench,bianchi2024negotiationarena}. It adds a connected account of one complete economy, from every offer through every traded monopoly to every bankruptcy.
